\documentclass[UTF8,a4paper,12pt]{ctexart}

\usepackage{amsmath,amssymb,bm}
\usepackage{booktabs}
\usepackage{geometry}
\usepackage{longtable}
\usepackage{array}
\usepackage{enumitem}
\usepackage{microtype}
\usepackage{hyperref}

\geometry{left=2.6cm,right=2.6cm,top=2.5cm,bottom=2.5cm}
\setlength{\parindent}{2em}
\setlength{\parskip}{0.35em}
\renewcommand{\arraystretch}{1.35}
\setlist{nosep}
\hypersetup{hidelinks}

\title{方法与实验部分}
\author{}
\date{}

\begin{document}

\maketitle

\setcounter{section}{2}

\section{提出的模型}

\subsection{站内疏散网络与共享设施容量耦合模型}

\subsubsection{问题定义}

本文研究大型多线换乘站中的整体应急疏散问题。与只根据空间距离为乘客分配出口的静态路径组织不同，本文将乘客的移动路径和楼梯、扶梯、换乘通道、闸机及出入口等有限通行设施的服务过程放在同一个疏散模型中描述。

在五线换乘站内，来自不同线路站台、站厅和换乘方向的客流可能先后或同时进入同一设施。某条路线在几何上能够到达出口，并不意味着该路线上的乘客能够在预计时刻获得设施服务并进入下游空间。因此，模型需要解决的不是单独的最短路径问题，而是：在多源客流持续进入的情况下，如何根据客流到达设施时的服务负荷组织可执行路径，并在实际容量约束下完成整体疏散。

\subsubsection{站内网络表示}

将龙阳路站的站台、站厅、换乘连接、垂向设施、闸机、出口前区域和出入口抽象为有向图
\[
G=(V,E),
\]
其中，节点集合 \(V\) 表示站台等待区、列车门区、站厅、换乘连接节点、楼梯或扶梯入口、闸机、出口前连接区及出入口；边集合 \(E\) 表示相邻空间或设施之间的有向移动关系。每条边 \(e\in E\) 包含长度 \(l_e\)、宽度或等效宽度 \(w_e\)、边类型、行进方向、所属区域和物理资源标识等属性。

对于具有独立通行能力的设施，建立边到共享资源的映射
\[
\rho:E\rightarrow R,
\]
其中 \(R\) 为有限通行设施资源集合。映射到同一资源 \(r\in R\) 的多条上游边共同使用该资源的服务能力。例如，多个上游节点连接到同一楼梯入口时，不能分别为每条边计算一份独立楼梯容量，而应在资源层面对进入请求进行统一分配。该处理使站台客流、站厅客流和跨线换乘客流在共享设施上的竞争能够被显式记录。

\subsubsection{设施服务容量与下游接收约束}

设资源 \(r\) 的服务率为 \(\mu_r\)，仿真时间步长为 \(\Delta t\)，上一时间步未使用的容量余量为 \(\xi_r(t)\)，则时间步 \(t\) 内资源能够释放的整数容量表示为
\[
C_r(t)=\left\lfloor \mu_r\Delta t+\xi_r(t)\right\rfloor .
\]

释放容量扣除后保留新的余量。若多个上游节点同时向资源 \(r\) 发送移动请求，所有请求先汇总到资源层，再按照整数容量进行分配。未被接受的请求不能被路径搜索结果直接视为已经通过，而应继续留在上游等待状态或资源队列中。

设施服务能力还受到下游空间接收条件的约束。设节点 \(v\) 的当前占用人数为 \(N_v(t)\)，在当前时间步被允许进入该节点的人数为 \(y_v(t)\)，节点接收上限为 \(\bar N_v\)，则需要满足
\[
N_v(t)+y_v(t)\leq \bar N_v.
\]

因此，路径层产生的是“希望进入某条边或某项设施”的移动请求，执行层还要根据资源剩余容量、目标节点占用和边内空间状态决定实际接受人数。路径可达性与实际通行量由此被区分开来。

\subsubsection{物理移动时间与速度--密度关系}

对于具有实际空间长度的移动边，根据边内客流密度计算物理行走速度。本文采用分段速度--密度关系：
\[
v(k)=
\begin{cases}
v_f, & k\leq k_f,\\
\max(v_f-ak,0), & k_f<k<k_j,\\
0, & k\geq k_j,
\end{cases}
\]
其中，\(k\) 为边内客流密度，\(v_f\) 为自由流速度，\(k_f\) 为自由流密度阈值，\(a\) 为速度下降斜率，\(k_j\) 为拥挤密度上限。边 \(e\) 在时刻 \(t\) 的物理行走时间为
\[
\tau_e(t)=\frac{l_e}{v(k_e(t))}.
\]

模型采用 \(v_f=1.427\ \mathrm{m/s}\)、\(k_f=0.2\ \mathrm{p/m^2}\)、\(a=0.3549\ \mathrm{m^2/(p\cdot s)}\)、\(k_j=4.0\ \mathrm{p/m^2}\)，以及 \(3.0\ \mathrm{p/m^2}\) 的高密度判定阈值。上述参数分别属于速度--密度关系的复现参数或模型判定参数，不作为龙阳路站实测值使用。楼梯、扶梯等设施的速度上限属于模型设定，除非有正式测量或运营资料支撑，不将其表述为车站官方数据。

需要特别区分的是，设施排队人数、边内占用人数和空间密度不是同一个状态变量。排队人数用于描述服务尚未释放的需求，边内占用和节点占用用于描述空间接收状态。将排队人数直接折算成下游密度会重复计算拥堵，因此本文在执行层分别记录资源队列和空间阻塞。

\subsection{面向设施到达时刻的队列预测与动态路径组织模型}

为使路径评价能够反映客流到达设施时的实际服务负荷，本文根据客流批次的预计到达时刻和已确定的设施服务事件，构造确定性的设施到达状态预测。该模块将客流当前可观测的设施状态转换为其预计到达设施时可能面对的服务状态。

\subsubsection{多源客流与自然到达批次}

客流不被简化为同一时刻从同一起点出发的总人数，而是由来源组和自然到达批次共同表示。一个批次记为
\[
b=(g,n,x,t_0,P,r,t_q),
\]
其中，\(g\) 为来源组，\(n\) 为批次人数，\(x\) 为当前节点，\(t_0\) 为该批次可以进行路径决策的时间，\(P\) 为当前路径，\(r\) 为已选择的共享设施资源，\(t_q\) 为进入资源队列的时间。来源组用于区分不同线路站台客流、站厅客流、列车到达客流和跨线换乘客流，并保留换乘方向及其在站内网络中的起点。

在同一仿真时间步内，批次按照预计到达时间、来源组、当前节点和批次编号的确定性顺序处理。批次完成路径选择后，其整数人数和预计到达时间立即登记为资源到达事件，使后续批次能够看到本轮已经作出的设施使用承诺。尚未形成明确路径的需求、浮点预留人数和未被接受的移动请求不进入确定性预测队列，避免对同一客流重复计数。

\subsubsection{设施到达时的队列预测}

设资源 \(r\) 的当前实际队列为 \(Q_r(t)\)，候选批次预计在时刻 \(\eta\) 进入该资源。预测队列由当前实际队列、尚未消耗资源入口容量的在途到达事件，以及本轮此前已经完成路径分配的整数批次共同构成。用队列演化算子表示为
\[
\widehat Q_r(\eta)=\mathcal{Q}\left(Q_r(t),\{(\eta_j,n_j)\mid \eta_j\leq\eta\},\mu_r\right),
\]
其中，\(\mathcal{Q}(\cdot)\) 按时间推进资源服务、扣减已释放容量，并在相应到达事件发生时加入批次人数。

若 \(\mu_r>0\)，候选批次到达资源 \(r\) 时因已有需求产生的预计等待时间为
\[
W_r(\eta)=\frac{\widehat Q_r(\eta)}{\mu_r}.
\]

该等待时间不是当前时刻队列长度的静态换算，而是面向候选批次未来到达时刻的队列状态估计。它使路径前段的行走时间、前序设施等待和后续共享设施服务负荷能够连续传递到后续决策。

\subsubsection{时间依赖路径评价}

对于批次 \(b\) 的候选路径 \(P=(e_1,e_2,\ldots,e_m)\)，设其在第 \(i\) 条边或设施上的预计进入时间为 \(\eta_i\)，则路径评价函数写为
\[
J(P\mid b)=\sum_{i=1}^{m}
\left[\tau_{e_i}(\eta_i)+W_{r_i}(\eta_i)+S_{r_i}(n_b)+W^{\mathrm{sp}}_{e_i}(\eta_i)\right],
\]
其中，\(\tau_{e_i}\) 为物理行走时间，\(W_{r_i}\) 为批次到达资源时的预计队列等待，\(S_{r_i}(n_b)\) 为当前批次占用资源的服务时间，\(W^{\mathrm{sp}}_{e_i}\) 为目标节点存储或下游空间接收不足造成的预计等待。本文将批次服务时间表示为
\[
S_{r_i}(n_b)=\frac{n_b}{2\mu_{r_i}}.
\]

每一段的行走时间和等待时间都会更新后续边的预计进入时刻，因此路径评价的是整条路径上的累计到达状态，而不是分别比较每条边在当前时刻的局部代价。

\subsubsection{动态改路条件}

动态组织并不意味着乘客在移动过程中可以随时改变方向。本文只对仍处于明确决策阶段、尚未进入在途执行队列的批次重新评估路径；已经进入在途执行的批次继续完成已经承诺的边。

对于允许重评估的批次，设保留原路径的剩余成本为 \(C_{\mathrm{stay}}\)，候选路径的成本为 \(C_{\mathrm{switch}}\)，相对收益为
\[
R=\frac{C_{\mathrm{stay}}-C_{\mathrm{switch}}}{C_{\mathrm{stay}}}.
\]

只有在候选路径有效、下一跳确实发生变化、批次仍处于允许决策阶段、\(C_{\mathrm{switch}}<C_{\mathrm{stay}}\)，并且 \(R>g_{\min}\) 时才接受改路。当前 Mode 4 使用的 \(g_{\min}=0.20\) 是本项目的场景防摆动设定，不将其表述为外部校准得到的通用参数。

\subsection{车站疏散效果综合评价方法}

为全面评价路径组织对整体疏散过程的影响，本文按照“整体疏散效率--共享设施瓶颈--来源组组织效果--跨尺度一致性”建立评价体系，不预先假定某一个出口均衡指标一定改善。

\subsubsection{整体疏散效率}

采用达到指定疏散比例的时间描述整体清空过程，包括 \(T_{95}\)、\(T_{99}\) 和 \(T_{100}\)。其中，\(T_{100}\) 为最后一名目标乘客完成疏散的时间，\(T_{95}\) 和 \(T_{99}\) 用于识别尾部客流的清空延误。

累计停滞人秒定义为
\[
I_{\mathrm{stat}}=\int_{0}^{T_{100}}N_{\mathrm{stat}}(t)\,\mathrm{d}t,
\]
其中 \(N_{\mathrm{stat}}(t)\) 表示时刻 \(t\) 处于停滞或等待状态的乘客数。结果同时报告人均停滞时间、平均移动时间、平均总疏散时间和总移动距离。资源队列人秒作为设施诊断指标单独报告，不与整体停滞人秒重复相加。

\subsubsection{关键设施拥堵和负荷分配}

对每个有限服务资源报告累计通过人数、峰值队列、队列人秒、首次排队时间、最后排队时间、利用率和最大预计等待。对关键楼梯、换乘通道和闸机的负荷分布采用 Jain 指数：
\[
J(x)=\frac{(\sum_{i=1}^{n}x_i)^2}{n\sum_{i=1}^{n}x_i^2}.
\]

关键设施 Jain 指数用于判断客流是否过度集中于少数共享设施；出口 Jain 指数仅作为辅助诊断，不将其改善预设为路径组织方法必然产生的效果。设施负荷是否改善，必须与队列人秒、峰值队列和整体尾部时间共同解释。

\subsubsection{线路、客流来源与换乘客流指标}

本文仍以全站整体疏散指标作为主要评价对象，同时按线路、站台、站厅和换乘来源报告 \(T_{95}\)、清空时间、设施通过量、出口分配和路径分配。换乘客流不是另一个独立的疏散问题，而是大型换乘站中解释共享设施竞争的重要客流来源。来源组结果用于回答：换乘客流与其他线路客流在哪些设施上汇聚，路径组织改变是否改变了这些客流的到达时段和等待分布。

\section{动态路径组织模型的求解方法}

\subsection{时间依赖路径模型的可执行化转换}

本文不存在需要进行鲁棒对偶化的半无限规划模型，因此将设施到达状态预测转换为可执行的时间依赖路径决策，并进一步说明其搜索状态、标签筛选和执行约束。

\subsubsection{时间依赖搜索状态}

对批次 \(b\)，路径搜索标签至少包含当前节点 \(v\)、累计到达时间 \(\eta\)、累计路径代价 \(g\)、预计资源状态、已使用路径节点和前驱信息，表示为
\[
L=(v,\eta,g,\mathcal{Q},P,\pi).
\]

当标签从节点 \(u\) 沿边 \(e=(u,v)\) 转移时，先根据边内密度计算物理行走时间，再查询资源 \(r=\rho(e)\) 在预计进入时刻的队列状态，得到新的到达时刻
\[
\eta'=\eta+\tau_e(\eta)+W_r(\eta)+S_r(n_b)+W^{\mathrm{sp}}_e(\eta),
\]
并将相应代价累加到 \(g\)。静态自由流距离只用于构造搜索下界，不能替代资源到达时刻的动态队列评价。

\subsubsection{标签筛选和可行性约束}

若两个标签到达同一节点，其中一个标签在累计时间、累计代价、资源等待和路径可行性方面均不优于另一个标签，则将其删除。对存在不同到达时刻或不同资源使用历史的标签予以保留，以避免把未来队列状态不同的路径错误合并为一个静态节点状态。

搜索过程中同时检查路径回环、目标节点接收能力、设施方向和资源服务条件。搜索得到的路径只是批次的候选移动组织方案，不能绕过后续执行层直接产生实际通行人数。

\subsubsection{对照方法与正式方法}

ImprovedAStar 使用当前正式代码中的路径代价和速度--密度关系，在给定网络状态下计算客流到出口的路径，不使用 AA 的未来资源到达事件索引、批次级预计队列和动态改路机制。它是计算对照方法，不是龙阳路站现实中已经存在的官方疏散方案。

AdaptiveQueueAwareAStar（以下简称 AA）使用上述时间依赖搜索，根据批次预计到达共享设施时的队列和空间接收状态评价候选路径，并在满足改路条件时更新尚未进入在途执行的批次。两种方法的差异限定在路径评价、未来队列使用和路径更新机制，不改变同一客流输入下的设施容量与物理执行规则。

\subsection{模型求解与疏散执行流程}

每个仿真时间步按照“状态释放--批次决策--容量执行--状态更新”的顺序推进：
\begin{enumerate}
    \item 释放已经到达的在途客流，更新节点占用、边内人数和实际资源队列；
    \item 根据自然到达时刻生成客流批次，并按预计到达时间、来源组和批次编号确定处理顺序；
    \item 对每个可决策批次调用相应路径方法，计算候选路径的累计预计到达时间和资源队列等待；
    \item 将本轮已经完成路径分配的整数批次登记为确定的资源到达事件，使后续批次使用更新后的未来队列；
    \item 汇总不同上游节点对同一共享资源的移动请求，按照统一的整数容量分配规则接受或拒绝请求；
    \item 对目标节点存储、边内密度和空间回溢条件进行二次检查，将未接受人数保留在上游等待状态；
    \item 记录整体疏散、资源队列、设施负荷、来源组流量和路径变更事件；
    \item 对仍处于允许决策阶段的批次进入下一轮重评估。
\end{enumerate}

ImprovedAStar 和 AA 完全共用这一执行流程。路径搜索不直接释放楼梯、扶梯、闸机或出口容量，只有执行层接受的移动请求才计入实际疏散过程。该设置保证两种方法之间的比较反映路径组织差异，而不是不同的物理执行条件。

\section{案例研究}

\subsection{关键设施到达负荷预测结果分析}

\subsubsection{案例站点与网络构建}

本文以龙阳路站为案例对象，模型网络包含 L2、L7、L16、L18 和 Maglev 五类线路及其站台、站厅、换乘连接、垂向设施、闸机和出入口。案例建模的重点不是仅复刻车站平面几何，而是识别不同来源客流共同使用的设施资源，并保证同一资源在路径层和执行层使用同一容量口径。

站点背景、线路关系和设施信息应由运营方或政府公开资料支撑。模型参数按照来源分为三类：

\begin{table}[htbp]
\centering
\caption{模型参数的来源分类}
\begin{tabular}{p{0.22\textwidth}p{0.32\textwidth}p{0.34\textwidth}}
\toprule
参数类别 & 当前用途 & 本文处理方式 \\
\midrule
文献复现参数 & 速度--密度关系及其基础阈值 & 引用原始文献，不称为龙阳路站实测值 \\
模型设定参数 & 设施速度上限、仿真时间步、部分容量表达和搜索阈值 & 明确写为本文模型设定 \\
场景假设参数 & 队列存储面积、缺少实测资料的局部宽度或容量 & 明确写为假设，并通过敏感性分析检验 \\
\bottomrule
\end{tabular}
\end{table}

\(\Delta t=1\ \mathrm{s}\) 为离散仿真步长；速度--密度参数、设施速度上限和 \(g_{\min}=0.20\) 等数值分别按照文献复现参数、模型设定参数或场景假设参数处理。默认面积、回退宽度和容量不作为官方实测值使用。

\subsubsection{队列预测和路径组织机制验证}

本节验证确定性设施到达状态预测和动态路径组织机制。实验围绕以下问题展开：客流批次真正到达共享设施时，模型是否已经把前序批次的承诺到达和服务消耗纳入队列判断；队列预测是否改变了路径评价；路径改变是否能够由设施负荷变化解释。

具体分析包括：
\begin{enumerate}
    \item 对关键楼梯、换乘通道和闸机，绘制实际队列、在途确定到达、本轮已分配到达和预测队列的时间演化；
    \item 对同一来源组在不同到达时刻的候选路径，分解物理行走时间、资源队列等待、当前批次服务时间和下游空间接收等待；
    \item 统计 AA 的路径重评估、满足改路条件的事件、实际获得移动容量的改路事件和无可行替代路径事件；
    \item 检查改路是否只发生在允许的决策阶段，是否存在重复计数、在途瞬时改路、路径回环或容量超发；
    \item 对比相同来源组和相同批次在 ImprovedAStar 与 AA 下的关键设施选择、到达时段和出口分配。
\end{enumerate}

本节的判断标准不是 AA 调用了多少次 A*，而是预测队列是否改变了批次对共享设施通行条件的判断，并且这种判断是否在执行层表现为可解释的设施负荷变化。

\subsection{模型方案性能分析}

\subsubsection{实验分组与共同条件}

本文在 Mode 1 和 Mode 4 两种负荷条件下分别运行 ImprovedAStar 与 AA。四组实验使用同一站内网络、同一来源组客流定义、同一速度--密度关系、同一设施服务率、同一下游接收约束和同一仿真时间步。除路径评价和路径更新机制外，不改变物理执行条件。

Mode 1 表示低负荷场景，采用各线路站台等待、站厅和换乘基础客流，不加入满载列车客流，场景输入总人数为 2,187 人。该场景用于检验方法在设施竞争较弱时的组织表现，并观察动态组织是否引入额外绕行或扰动。

Mode 4 表示高负荷场景，在基础站内客流上叠加双向满载列车客流，场景输入总人数为 17,905 人，并启用 L2 列车来源按区域划分的客流表示。该场景用于放大多线路客流在楼梯、换乘通道、闸机和出入口前的竞争，检验设施到达负荷预测对整体尾部疏散和关键设施负荷的影响。

2,187 人和 17,905 人为场景输入，不是实测客流结果；其生成规则和场景含义在参数表中说明。低负荷和高负荷不是两个不同的研究问题，而是同一整体疏散问题的两种负荷条件。

\subsubsection{整体疏散效率比较}

首先比较 \(T_{95}\)、\(T_{99}\)、\(T_{100}\)、累计停滞人秒、人均停滞时间、平均移动时间、平均总疏散时间和总移动距离。该组实验回答的是：在相同设施容量和客流输入下，面向设施到达时刻的路径组织是否改变了整体疏散过程，尤其是高负荷场景中的尾部清空和停滞等待。

整体疏散性能的比较指标和结果记录格式如表所示：

\begin{table}[htbp]
\centering
\caption{两种负荷场景下的整体疏散性能比较}
\small
\begin{tabular}{llrrrrrrr}
\toprule
场景 & 方法 & 人数 & \(T_{95}\) & \(T_{99}\) & \(T_{100}\) & 停滞人秒 & 人均停滞 & 总移动距离 \\
\midrule
Mode 1 & ImprovedAStar & 2187 & 待填 & 待填 & 待填 & 待填 & 待填 & 待填 \\
Mode 1 & AA & 2187 & 待填 & 待填 & 待填 & 待填 & 待填 & 待填 \\
Mode 4 & ImprovedAStar & 17905 & 待填 & 待填 & 待填 & 待填 & 待填 & 待填 \\
Mode 4 & AA & 17905 & 待填 & 待填 & 待填 & 待填 & 待填 & 待填 \\
\bottomrule
\end{tabular}
\end{table}

结果分析不预设“AA 必然更优”。应先说明低负荷和高负荷下整体疏散过程的主要特征，再判断路径组织对尾部时间、等待和行走距离的影响是否具有一致性；若方法在不同指标之间存在权衡，应保留这种权衡，而不是只报告有利于某一方法的指标。

\subsubsection{关键设施负荷和客流分配比较}

进一步比较各共享设施的峰值队列、队列人秒、最大利用率、最后排队时间和关键设施 Jain 指数，并报告各线路、站台、站厅和换乘来源在关键设施上的通过量和等待分布。

该分析用于判断整体指标变化的站内原因：整体尾部延误是否主要由少数关键设施的持续排队造成；AA 是否减少了关键设施的集中使用，还是只是将队列转移到其他设施；路径组织是否通过增加有限的行走距离换取较低的设施等待。

换乘客流在此作为整体疏散中的重要来源组进行分析，而不是把实验目标改成“换乘客流单独疏散”。需要重点观察跨线换乘来源与站台、列车来源在共享楼梯、换乘通道和闸机上的共同使用关系，以及这种关系对全站尾部疏散的贡献。

\subsection{路径组织与乘客疏散综合效果分析}

本节将路径组织、设施执行、整体疏散和微观行人运动结合起来，分析路径方案对乘客疏散效果的综合影响。

\subsubsection{站内瓶颈与来源组机制}

按照队列人秒和峰值队列识别主要瓶颈设施，比较两种方法在这些设施上的客流到达时段、服务释放和上游等待，并结合线路和来源组路径分配分析客流汇聚过程。结果至少回答以下问题：
\begin{enumerate}
    \item 整体尾部延误是否由少数楼梯、换乘通道或闸机的持续排队造成；
    \item 换乘客流与其他线路客流在哪些共享设施上发生汇聚；
    \item AA 的路径变化是否真正降低了关键设施的集中到达，还是仅将排队转移到其他设施；
    \item 设施等待、空间阻塞和额外行走距离分别对整体疏散时间产生了什么作用。
\end{enumerate}

建议使用资源队列曲线、设施负荷表、来源组路径分配图、线路清空时间和关键设施到达热区图进行说明。出口利用和出口 Jain 指数只作为辅助结果，不把出口均衡预设为方法必然效果。

\subsubsection{速度--密度阈值敏感性分析}

当前敏感性分析围绕速度--密度模型中的高密度判定阈值进行扰动，例如比较约 2.0、3.0 和 3.9 \(\mathrm{p/m^2}\) 的设定，观察 \(T_{95}\)、\(T_{100}\)、停滞人秒、关键设施队列和方法间差异是否保持。

该实验用于检验结论对密度判定参数的敏感程度，属于参数稳健性分析。若后续要声称方法具有明确的需求--容量适用边界，还需要进一步改变客流规模、共享设施服务率或客流到达重叠程度，观察路径组织收益从明显到有限的变化区间。仅凭当前阈值扰动不能直接推出完整运行边界。

\subsubsection{Pathfinder 交叉验证}

Pathfinder 使用宏观疏散模型实际执行并导出的完整路径分配结果创建 Occupant Group 和 Behavior。导出的路线人数采用实际执行的整数人数，来源组内路线比例用于配置检查和结果报告，不重新生成一套与宏观模型不同的手工路线。

Pathfinder 不作为真实车站观测的替代物，也不作为宏观模型的绝对真值，而是从微观行人运动尺度检验以下趋势：
\begin{itemize}
    \item 总体疏散曲线和尾部清空趋势；
    \item 关键换乘设施、楼梯和闸机附近的拥堵空间分布；
    \item 不同路径组织方法的出口利用和路径使用差异；
    \item 宏观资源队列识别出的瓶颈位置是否在微观轨迹中出现相近的拥堵趋势。
\end{itemize}

若宏观模型和 Pathfinder 在趋势上保持一致，说明路径组织和瓶颈识别具有跨尺度一致性；若出现差异，则从微观行人行为、空间几何和容量映射差异解释，不能直接将 Pathfinder 视为宏观模型的绝对验证真值。

\subsubsection{结果解释顺序}

实验结果分析按照以下顺序展开：
\begin{enumerate}
    \item 先验证设施到达负荷预测是否能够识别未来队列，并确认路径变化发生在允许的决策阶段；
    \item 再比较 Mode 1 和 Mode 4 下两种模型方案的整体疏散性能；
    \item 接着结合关键设施、线路和换乘来源组解释整体差异从何而来；
    \item 然后报告速度--密度阈值敏感性，说明参数扰动下结论是否保持；
    \item 最后用 Pathfinder 从微观行人运动尺度检验宏观疏散趋势和瓶颈位置。
\end{enumerate}

上述实验依次从设施到达负荷、模型方案性能、关键设施与来源组机制、参数稳健性和微观行人运动五个层面评价路径组织方法。

\end{document}
